\begin{table}[t]
\centering
\scriptsize
\setlength{\tabcolsep}{2pt}
\resizebox{\columnwidth}{!}{%
\begin{tabular}{@{}llrrrrrrr@{}}
\toprule
\textbf{Task} & \textbf{metric} & \textbf{base} & \textbf{k8} & \textbf{k12} & \textbf{k14-p} & \textbf{k14-f} & \textbf{frozen} & \textbf{random} \\
\midrule
BoolQ & raw & .815 & .588 & .610 & .606 & .638 & .592 & .614 \\
BoolQ & norm. & .753 & .620 & .657 & .682 & .689 & .700 & .619 \\
\addlinespace[1pt]
CSQA & raw & .665 & .442 & .508 & .506 & .499 & .453 & .450 \\
CSQA & norm. & .652 & .405 & .471 & .479 & .476 & .467 & .405 \\
\addlinespace[1pt]
SIQA & raw & .502 & .400 & .415 & .441 & .434 & .414 & .416 \\
SIQA & norm. & .547 & .431 & .460 & .474 & .474 & .451 & .443 \\
\bottomrule
\end{tabular}%
}
\caption{\textbf{Raw versus continuation-length-normalized accuracy.} The normalized score divides candidate log-likelihood by continuation character length before selecting an answer. The main results report raw accuracy for BoolQ, CommonsenseQA, and SocialIQA. Normalization can materially change BoolQ and SIQA, so both views are reported here. k14-p/k14-f denote keep14 at 153.5k/200k; frozen-front and random init are the other 200k operating points.}
\label{tab:app-metric-sensitivity}
\end{table}
